\documentclass{article}
\usepackage{graphicx} % Required for inserting images

% Added stuff
\usepackage{amsmath}
\usepackage{amssymb}
\newcommand{\N}{\mathbb{N}}
\newcommand{\Z}{\mathbb{Z}}
\newcommand{\Q}{\mathbb{Q}}
\newcommand{\R}{\mathbb{R}}
\newcommand{\C}{\mathbb{C}}
\newcommand{\bigp}[1]{\left(#1\right)}
\newcommand{\bigb}[1]{\left[#1\right]}
\newcommand{\bigabs}[1]{\left|#1\right|}
\newcommand{\eps}{\varepsilon}
\renewcommand{\emptyset}{\varnothing}

\title{p-adic Analysis 5}
\author{Siddharth K}
\date{May 2026}

\begin{document}

\maketitle

\section*{Problem 2}
Let $x = 9/16$ and $y = 9/4$. The $p$-adic logarithm is
$$\log(1 + z) = \sum_{n = 1}^\infty (-1)^{n + 1} \frac{z^n}{n},$$
so
$$\sum_{n = 1}^\infty (-1)^n \frac{3^{2n}}{n4^{2n}} = \sum_{n = 1}^\infty (-1)^n \frac{x^n}{n} = -\log(1 + x).$$
Since $x = 9/16$,
$$1 + x = \frac{25}{16} = \bigp{\frac{5}{4}}^2,$$
$$-\log(1 + x) = -\log\bigp{\bigp{\frac{5}{4}}^2}.$$
Also,
$$1 - y = -\frac 54,$$
so
$$(1 - y)^2 = \frac{25}{16} = 1 + x.$$
Therefore,
$$-\log(1 + x) = -\log((1 - y)^2) = -2\log(1 - y)$$
$$= -2\sum_{n = 1}^\infty \frac{3^{2n}}{n4^n}.$$
This highlights that $\log(-5/4) = \log(5/4)$ because $\log(-1) = 0$.

\section*{Problem 4}
For $n \ge 2$,
$$v_p\bigp{\frac{p^n}{n!}} \ge n - v_p(n!) \ge 2,$$
so
$$e^p \equiv 1 + p \pmod{p^2}.$$
For contradiction, assume that $x^p \equiv e^p \pmod{p^2}$ for some $x \in \Z_p$. Since
$$e^p \equiv 1 + p \equiv 1 \pmod{p},$$
we have that
$$x^p \equiv 1 \pmod{p}.$$
Also,
$$x^p \equiv x \pmod{p},$$
$$x \equiv 1 \pmod{p}.$$
So we can write $x = 1 + pa$ with $a \in \Z_p$. So
$$x^p = 1 + p(pa) + \binom{p}{2}p^2a^2 + \cdots \equiv 1 \pmod{p^2}.$$
But we already found $e^p \equiv x^p \equiv 1 + p \pmod{p^2}$, which is a contradiction.

\section*{Problem 5}
Assume $x \equiv y \pmod{p^N}$, so that $x$ and $y$ have the same first $N$ digits. That means $f(x)$ and $f(y)$ have the same digit contributions in positions
$$0^2, 1^2, \ldots, (N - 1)^2,$$
so
$$f(x) \equiv f(y) \pmod{p^{N^2}}.$$
Therefore,
$$|x - y|_p \le p^{-N} \implies |f(x) - f(y)|_p \le p^{-N^2},$$
which proves continuity.

For differentiability, let $h \in p^N\Z_p$. Then $x$ and $x + h$ have the same first $N$ digits. Therefore,
$$f(x + h) - f(x) \in p^{N^2}\Z_p,$$
so
$$\bigabs{\frac{f(x + h) - f(x)}{h}}_p \le p^{-N^2 + N} \to 0,$$
so $f'(x) = 0$.

\section*{Problem 12}
For any $x \in B$ and $m \ge 0$, we have that
$$f(x + mk) = f(x).$$
Now, fix $x \in B$. For all $r \ge 0$, we can take $m = p^r$, so that
$$f(x + p^rk) = f(x),$$
$$\frac{f(x + p^rk) - f(x)}{p^rk} = 0.$$
As $r \to \infty$, $|p^rk|_p \to 0$, so $p^rk \to 0$. So the derivative of $f$ at $x$ can be written as
$$f'(x) = \lim_{r \to \infty} = \lim_{r \to \infty} \frac{f(x + p^rk) - f(x)}{p^rk} = 0.$$
So $f'(x) = 0$ for all $x \in B$. Since $f$ can be written as a power series, all coefficients of terms of degree $\ge 1$ must be $0$, and $f$ is constant.

\section*{Problem 14}
Take $B = \Z_p$ and $f(x) = x^2$. This $f$ is analytic on $B$. Now, choose $a = 0$, $b = 1$, $n = -1$. Then
$$|f(a)|_p = 0, \quad |f(b)|_p = |1|_p = 1.$$
So it is true that
$$|f(a)|_p \le p^{-1} \le |f(b)|_b.$$
IVT in $p$-adics would imply that there exists $c \in \Z_p$ where
$$|f(c)|_p = p^{-1},$$
but
$$|f(c)|_p = |c^2|_p = |c|_p^2 \in \{0, 1, p^{-2}, p^{-4}, \ldots\},$$
where none of the options are $p^{-1}$. So IVT fails, even for analytic functions.

\section*{Problem 16}
Take $B = \Z_p$ and $f(x) = x^p$, which is analytic. Choose $a = 0$, $b = 1$. So
$$\frac{f(b) - f(a)}{b - a} = 1.$$
Assuming MVT in $p$-adics would imply that there is $x \in \Z_p$ such that $f'(x) = 1$. But $f'(x) = px^{p - 1}$, so
$$x^{p - 1} = \frac{1}{p}.$$
Since $|x|_p \le 1$,
$$|x^{p - 1}|_p = |x|_p^{p - 1} \le 1,$$
but
$$\bigabs{\frac{1}{p}}_p = p > 1.$$
This is a contradiction, so no such $x$ exists and MVT is not true for $p$-adics.

\end{document}